% The class options come from tools/build.sh: sigconf,anonymous,review for
% the review version (the default), sigconf otherwise.
\providecommand\PaperClassOptions{sigconf,anonymous,review}
\expandafter\documentclass\expandafter[\PaperClassOptions]{acmart}

\usepackage{grants}
\usepackage{booktabs}
\usepackage{tabularx}
\usepackage{multirow}
\usepackage{xltabular}

% Two narrow columns full of code identifiers: let TeX stretch a line a
% little rather than overrun the margin.
\setlength\emergencystretch{1.5em}

\setcopyright{none}
\settopmatter{printacmref=false}
\renewcommand\footnotetextcopyrightpermission[1]{}
\acmConference{}{}{}
\acmYear{}
\acmDOI{}
\acmISBN{}

% Theorems 1--4, numbered on their own so the paper and the slides agree.
\AtEndPreamble{\theoremstyle{acmplain}\newtheorem{mainthm}{Theorem}}

\begin{document}

\title{One Policy, Two Boundaries: Verified Capability Enforcement for Agents Across WebAssembly and Smart Accounts}

\author{Nido contributors}
\affiliation{%
  \institution{Nido}
  \country{}}

\begin{abstract}
AI agents now hold real authority: they read and write files, call web
services and sign blockchain transactions. The controls in use check that
authority at the wrong grain for Soroban smart accounts. A custody policy
engine such as Dfns applies its amount and recipient rules only to plain
transfers, and a process sandbox confines paths, processes and hosts but sees
a call's arguments, at most, as text. We write an agent's authority as a
\emph{grant}: a list of caveats over actions, which its holder can only
extend before handing a copy to a sub-agent, while issuing, revoking, key
rotation and budget refills stay in an admin document that agents never hold.
One evaluator, written in a restricted subset of Rust, decides grants at two
boundaries. In the Soroban smart account it runs as a CAP-0071 delegated
signer, which the host hands every context of the account's authorization, so
a call the account authorizes cannot hide a disallowed one in its tree. In the
agent's host it runs as a WebAssembly component in front of files, the network
and a key keeper. We prove that the evaluator's Charon and Aeneas translation
to Lean computes a readable model in four action vocabularies and for budgets,
so totality, attenuation, delegation chaining and budget safety under
oversubscribed allowances hold \translationtrust; a step bound is proven for a
counted model and checked against metered costs. In Soroban simulation all 45
escalation attempts are refused with state unchanged, as are all 31 at the
host and the 44 that run on a local Stellar network through the relay, where
sub-agents racing to overdraw their root are charged within its allowance;
the evaluator's metered work tracks the model's step count at about 99 CPU
instructions per step; an agent transaction costs 2.2M to 5.0M of Soroban's
400M per-transaction instructions in simulation, and 9 to 12\% more on the
network; and the relay confirms one transaction per channel per ledger. Every
claim carries its status:
proven, tested, measured, assumed or unproven.
\end{abstract}

\keywords{capabilities, attenuation, AI agents, smart accounts, Soroban,
WebAssembly, formal verification, Lean, Aeneas}

\maketitle

\section{Introduction}\label{sec:intro}

An agent that pays invoices, renews cloud resources or trades for a user needs
authority that is real and bounded. It must be able to read the files it works
on, reach the services it uses and sign transactions, but only those the user
meant. For an agent acting through a Soroban smart account, the controls in
use check that authority at the wrong grain.

\emph{Custody policy engines} guard the signing key. Dfns, for example,
evaluates rules such as transaction amount limits and recipient allow-lists
before it signs~\cite{dfnsPolicies}. Those rules ``only evaluate plain value
transfers'': when the amount or recipient cannot be inferred, as for a contract
call or a raw signature request, they ``fail closed and
trigger''~\cite{dfnsPolicies}, and the engine does not ``introspect raw
serialized signature payloads''~\cite{dfnsSigningPolicies}. Other engines
decode contract calls and write rules over methods and arguments
(\S\ref{sec:related}), but the ones whose documentation we read describe
EVM-style calls, and we found none that evaluates a Soroban authorization
tree, in which one signed call nests the calls it makes for the account.

\emph{Process sandboxes} guard the machine. nono confines an agent process with
Landlock on Linux and Seatbelt on macOS: it restricts file system paths,
process execution and network destinations, down to a domain, a URL path or an
HTTP method and path, and it matches command arguments against exact, prefix or
contains patterns~\cite{nono}. It does not read what a call carries, such as
the amount and recipient inside a transaction that a remote key signs, and the
smart account that finally executes the transaction cannot see the sandbox's
rules at all.

We write an agent's authority as a \emph{grant}: a list of caveats, each a
predicate over actions, where an action is allowed only if every caveat holds.
A grant's holder can do two things with it: act within it, and hand a sub-agent
a copy with more caveats added. Issuing, revoking, rotating a holder's key and
refilling a budget are the operations of a separate \emph{admin document} that
agents never hold. One evaluator decides every check, at two boundaries:
\begin{itemize}
\item In the \textbf{Soroban smart account}, the evaluator runs inside a
  delegated signer. Through CAP-0071 delegated signing~\cite{cap71,cap7101} the
  host hands it every context of the account's authorization, so a call the
  account authorizes cannot hide a disallowed one in its authorization tree.
\item In the \textbf{agent's host}, the same evaluator, compiled to the same
  WebAssembly target, runs as a component behind the host that performs the
  agent's file, network and signing requests. Before its key keeper signs a
  transaction or an authorization entry, the host checks every context the
  signature authorizes against the agent's chain grant (\S\ref{sec:host}).
\end{itemize}
A relay that submits transactions for agents holding no keys runs the same
chain check before a custody key signs and pays for them. Figure~\ref{fig:arch}
shows the three deployments.

\begin{figure*}[t]
\centering
\resizebox{0.8\textwidth}{!}{\input{architecture}}
\caption{One evaluator at every boundary an agent's actions cross. Each
boundary checks the agent's grant with the same evaluator before acting; the
admin document that changes grants is the owner's alone.}
\label{fig:arch}
\Description{An agent sends tool calls to the agent host, relay requests to
the relay, and holder signatures to the smart account; each boundary runs the
grant evaluator before reaching files, the network and the key keeper, the
custody key and channel accounts, or the contracts the account calls.}
\end{figure*}

The adversary is an agent the user already authorized, trying to widen its own
authority (\S\ref{sec:threat}): by editing its grant, by hiding a disallowed
call inside an allowed one, or by working around delegation through deployed
contracts, approvals or a sibling's allowance.

\paragraph{Contributions.}
\begin{enumerate}
\item A grant model with four action vocabularies (chain calls, files,
  network, signing) that share one evaluator, delegation by attenuation only,
  and budgets as allowances charged up the delegation tree, which stay safe
  when children's allowances oversubscribe their parent's
  (\S\ref{sec:grants}).
\item A proof from the deployed Rust back to a readable Lean model
  (\S\ref{sec:proofs}). Charon and Aeneas translate the evaluator to
  Lean~\cite{charon,aeneas,lean4}, and we prove that the translation computes
  the model for every input in every vocabulary and for the budget functions.
  Totality, attenuation, chaining and budget safety therefore hold
  \translationtrust. Attenuation and chaining are immediate in the model; what
  the proof adds is their carry-over to the code and the budget invariant
  under oversubscription. A step bound is proven for a counted model of the
  evaluator, not for the code.
\item Two enforcement boundaries built on one evaluator, with no change to the
  smart account or to OpenZeppelin's account library: a CAP-0071 delegated
  signer for Soroban (\S\ref{sec:soroban}) and a WebAssembly component behind
  an agent host (\S\ref{sec:host}), plus a relay that checks before a custody
  key signs (\S\ref{sec:relay}).
\item A separation of what is proven from what is tested, measured, assumed
  or unproven, with the trusted wiring and its known gaps stated
  (\S\ref{sec:tcb}), and an evaluation (\S\ref{sec:eval}): an attack suite of
  76 escalation attempts in Soroban simulation and at the host, all refused
  with state unchanged, the Soroban attempts again on a local Stellar network
  through the relay, with its throughput and a race of sibling spends, and a
  cross-check of the proven step bound against metered Soroban costs.
\end{enumerate}

\paragraph{Claim status.} A tag after a claim gives its status, recorded in a
claim ledger that the build checks against the repository: \claimtag{proven}
by named Lean theorems, \claimtag{tested} by a named test,
\claimtag{measured} by a recorded result file, \claimtag{assumed} as a stated
assumption, or \claimtag{unproven}. Appendix~\ref{app:claims} lists every
claim with its evidence. Grants control access, not information flow: a grant
decides which actions an agent may take, not where the data it reads goes next
through other allowed actions (\S\ref{sec:limits}).

\section{Background}\label{sec:background}

\paragraph{Soroban authorization.} A Soroban contract asks an address to
authorize the current call with \code{require\_auth}, which requires ``address
to have authorized call of the current contract function with all the
arguments it has been called with''~\cite{cap46auth}. A transaction carries
authorization entries. Each holds credentials and ``a single tree of contract
calls'' (\code{SorobanAuthorizedInvocation})~\cite{cap46auth}, and the host
matches every \code{require\_auth} during execution to a node of a signed
tree, so a nested call that the signed tree leaves out fails. A contract that
implements \code{\_\_check\_auth} is a custom account. When an authorization
names it, the host calls its \code{\_\_check\_auth} with a signature payload,
the signature, and ``a list of the contract invocations that are being
authorized by these signatures''~\cite{stellarAuthDocs}. Each such
\emph{context} is either a contract call (contract, function, arguments) or a
contract-creation host function; a \code{\_\_check\_auth} that returns an
error or traps fails the authorization~\cite{cap46auth}.

\paragraph{Delegated signing (CAP-0071).} Protocol 27 adds the host function
\code{delegate\_account\_auth}, which delegates ``the authentication logic from
an \code{Address}es \code{\_\_check\_auth} function to a different address'',
and credentials that carry the delegates' signatures inside the account's own
entry~\cite{cap71,cap7101}. The delegate's signature payload and authorization
context ``are equivalent to the signature payload and the context of the
`top-level' address''~\cite{cap7101}. So a contract delegate's
\code{\_\_check\_auth} receives the delegating account's payload and every
context of the account's authorization\claim{cap71-contexts}.

\paragraph{Smart accounts.} Nido's accounts are Perch accounts built on
OpenZeppelin's Stellar smart-account library, which ``organizes authorization
as a composition of three core elements: context rules, signers, and
policies''~\cite{ozStellarAccounts}. A context rule scopes a set of signers
and policy contracts to a kind of context, and the submitter selects one rule
per context. A Perch account takes all of its authorization data from one
policy document, whose only write path is the account's \code{apply\_doc}
function; a document can declare a signer as an address, which the account
authenticates through \code{delegate\_account\_auth}.

\paragraph{WebAssembly components.} WebAssembly modules execute in a linear
memory ``disjoint from code space, the execution stack, and the engine's data
structures''~\cite{wasm}. The component model types a module's imports and
exports in WIT, and ``the only ways a component can interact with anything
outside itself are by having its exports called, or by calling its
imports''~\cite{componentModelDocs,componentModel}. Wasmtime runs
components~\cite{wasmtime}.

\paragraph{Charon and Aeneas.} Charon extracts a Rust crate into LLBC, a clean
intermediate language~\cite{charon}. Aeneas translates LLBC into a pure
functional program in a proof assistant's language, Lean among
them~\cite{aeneas,lean4}. A translated function returns in an error monad:
\code{ok v} for a value, or a failure for a panic, an arithmetic overflow or an
out-of-bounds index, or divergence. A theorem that a translated function
returns \code{ok} for every input therefore says the Rust function neither
fails nor diverges, under the trust in the translation stated in
\S\ref{sec:tcb}.

\section{Threat Model}\label{sec:threat}

\paragraph{Principals.} The \emph{owner} holds the account's admin key (in Nido,
a passkey) and through it the admin document: the owner issues root grants,
revokes them, rotates holder keys and refills budgets. An \emph{agent} holds a
grant, identified at the Soroban boundary by an ed25519 holder key and at the
agent host and the relay by the identity the host or relay assigned it. The
agent host, the relay and the key keeper are operated for the owner.

\paragraph{Adversary.} The adversary is an agent the owner already authorized,
or anything that controls one, such as injected instructions in the data the
agent reads. It may choose any tool call, any transaction and any
authorization tree; sign anything with its holder key; delegate to sub-agents
it controls; and call any contract its grant reaches. Its goal is authority
beyond its grant:
\begin{enumerate}
\item \textbf{Self-escalation}: edit its own grant, directly or through the
  admin functions, the account's rules, or its key.
\item \textbf{Hidden calls}: make a call its grant allows which, inside the
  same Soroban authorization tree, makes a call its grant does not allow. The
  tree holds the calls that need the account's authorization; a call's effects
  that need none are outside it (below).
\item \textbf{Delegation workarounds}: escape its limits by deploying a
  contract, creating an account, handing authority to another address
  (an \code{approve} followed by \code{transfer\_from}), delegating from an
  ancestor's grant, or spending through siblings whose allowances together
  exceed their parent's.
\end{enumerate}

\paragraph{Goal.} At the Soroban boundary, every context that needs the
account's authorization on an agent's behalf is allowed by that agent's grant,
and the total spent of the budget asset at or under any grant stays within its
allowance. At the agent host, every file and network action the host performs
for an agent is allowed by the agent's grant, and every signature by its
signing caveats and, for a transaction or authorization entry, by its chain
caveats on every context the signature authorizes. At the relay, every context of
a relayed call is allowed by the agent's chain grant before the custody key
signs. Only the admin document changes a grant.

\paragraph{Assumptions.} The agent host is the agent's only path to files,
the network and keys\claim{sole-path}, and a deployment binds each agent's
tool-call channel to the agent id the host issued it\claim{agent-identity}. No
contract a grant allows can move the account's assets without the account's
authorization, through a standing token allowance or funds the account
deposited with it\claim{no-standing-allowance}; such effects need no context
of the account's authorization, so no caveat or budget sees them. The account's
document routes holder keys only through the grant signer and names its own
instance\claim{document-requirements}. The key keeper signs only what the host
or relay passes it\claim{keeper-trusted}. The Soroban host behaves as
specified, ed25519 signatures cannot be forged, and SHA-256 is collision
resistant. The compilers, translators and runtimes in \S\ref{sec:tcb} are
trusted.

\paragraph{Out of scope.} Information flow: an agent allowed to read a file
and to reach a host can send the file's contents there. The owner's choice of
grants: a grant that is too broad is enforced as written. Compromise of the
owner's admin key, the host, the relay or the key keeper. Confinement of the
agent process itself, which a process sandbox does and the host does not
(\S\ref{sec:host}). The repository's other relay, \file{infra/relayer}, which
fee-sponsors any contract call with valid on-chain authorization and runs no
grant check; accounts that use a grant signer are still checked on chain
whichever relay submits. Denial of service and side channels.

\section{Grants}\label{sec:grants}

\subsection{The model}

A grant over an action type $\alpha$ is a list of caveats, each a predicate on
$\alpha$. Figure~\ref{fig:model} is the readable Lean specification
(\file{formal/AttenuatingGrants/Grant.lean}), quoted from the repository. An
action is allowed when every caveat holds, so the empty grant allows every
action. A batch of actions, which is how every context of a Soroban
authorization tree is checked, is allowed when every action in it is. One
grant is narrower than another when every action the first allows, the second
allows too, and attenuation adds caveats.

\begin{figure}[t]
\begin{lstlisting}[language=Lean]
structure Grant (α : Type) where
  caveats : List (α → Bool)

def Grant.allows (g : Grant α) (a : α) : Bool :=
  g.caveats.all (· a)

def Grant.allowsAll (g : Grant α) (actions : List α) : Bool :=
  actions.all g.allows

def Grant.narrower (g h : Grant α) : Prop :=
  ∀ a, g.allows a = true → h.allows a = true

def Grant.attenuate (g : Grant α) (c : α → Bool) : Grant α :=
  { caveats := c :: g.caveats }
\end{lstlisting}
\caption{The readable grant model.}\label{fig:model}
\Description{Lean definitions of Grant, allows, allowsAll, narrower and attenuate.}
\end{figure}

\paragraph{Two documents.} The grant signer and the agent host split the
operations on grants between two interfaces. The \emph{admin document}
issues a root grant, revokes a grant (and with it every grant delegated from
it), rotates a holder's key and refills a budget. The \emph{agent interface}
acts under a grant and delegates: \code{delegate(add)} creates a child grant
whose caveats are the caller's followed by \code{add}, so the child allows
only what the caller allows, whatever \code{add} says. No agent operation
reads or edits a grant. At the relay, grants are a configuration file that
the admin side writes; agents only send requests.

\subsection{Four vocabularies}

Four action vocabularies plug into the model (Table~\ref{tab:vocab}). Each
defines an action type, a caveat type and a denotation from caveats to
predicates. Identifiers (contract ids, function names, wasm hashes, hosts,
key ids, network ids, path components) are byte strings compared by exact
equality. The evaluator normalizes nothing: lower-casing hosts and resolving
\code{..}, symlinks and case happen in the host before an action exists
(\S\ref{sec:host}).

\begin{table}[t]
\caption{The four action vocabularies.}\label{tab:vocab}
\small
\begin{tabularx}{\linewidth}{@{}l>{\raggedright\arraybackslash}X>{\raggedright\arraybackslash}X@{}}
\toprule
& \textbf{Action} & \textbf{Caveats} \\
\midrule
Chain & \code{invoke} contract function args, or \code{createContract} wasmHash
  & \code{contractIn}, \code{functionIn}, \code{argI128AtMost} $i$ $n$, \code{argEq} $i$ $v$ \\
Files & path (a list of components), mode $\in\{$read, write, execute$\}$
  & \code{pathUnder} prefix, \code{modeIn} \\
Network & host, port, scheme $\in\{$http, https$\}$
  & \code{hostIn}, \code{portIn}, \code{schemeIn} \\
Signing & key, network, payload $\in\{$transaction, authEntry, message$\}$
  & \code{keyIn}, \code{networkIn}, \code{payloadIn} \\
\bottomrule
\end{tabularx}
\end{table}

\paragraph{Chain.} Arguments follow Soroban's values: \code{i128},
\code{u64}, \code{bool}, \code{bytes}, \code{address} and \code{symbol} are
modelled, and an argument of any other type is \code{other} carrying its full
XDR encoding, whose leading discriminant names its type. Values of different
constructors are never equal, so an \code{argEq} caveat holds only for an
argument of its value's constructor\claim{arg-kind}, and an unmodelled
argument never matches a \code{bytes} caveat, whatever its
encoding\claim{other-not-bytes}. Every chain caveat is false on
\code{createContract}, so any chain grant with at least one caveat denies
contract deployment\claim{create-contract-denied}.

\paragraph{Caveats apply to every context alike.} A chain grant is checked
against every context of an authorization tree, and each caveat reads the same
fields of each context: \code{argI128AtMost i n} and \code{argEq i v} read
argument $i$ of every context, whatever its contract and
function\claim{positional-caveats}. Two consequences follow. A grant with
\code{contractIn [R, T]} and \code{functionIn [swap, transfer]} allows the
cross product, \code{R.transfer} and \code{T.swap} included. And a tree whose
calls carry amounts at different positions cannot have each amount capped: a
router's \code{forward} takes its amount as argument 3 and the token
\code{transfer} it nests as argument 2, so one \code{argI128AtMost} caps the
wrong argument of one of them. The evaluation's router workload therefore
bounds amounts by its allowance alone (\S\ref{sec:eval}). A caveat that
applies only to contexts of one contract and function would close this; the
generic model takes any predicate, so Theorems~\ref{thm:atten}
and~\ref{thm:chain} would be unchanged, but the vocabulary does not have one.

\paragraph{Files.} Paths are lists of components, so a grant under
\code{/home/a} never covers \code{/home/ab}\claim{path-components}. A
read-only grant is \code{modeIn [read]}.

\paragraph{Network and signing.} A network action is the host name the agent
requested, its port and its scheme; a signing action is a key id, a network
passphrase and a payload kind. Each caveat is membership of one field in a
list.

\subsection{Budgets}\label{sec:budgets}

A grant may carry an allowance of one asset. Budgets are enforced at the
Soroban boundary, where the grant signer meters transfers of one Stellar Asset
Contract per root grant; the agent host and the relay do not meter.
Allowances live on the delegation tree: the ledger holds one node per grant with a parent link, an
allowance and \code{spent}, the total of every successful spend at or under
the node. A spend at node $i$ walks from $i$ to its root and fails at the
first node whose \code{spent} plus the amount would exceed its allowance;
otherwise it charges every node on the walk. Either the spend succeeds and is
charged to the spender and every ancestor, or it fails and changes nothing.

Allowances may oversubscribe: children's allowances may sum to more than their
parent's, and delegation accepts any allowance. Because every ancestor is
charged in the same step, spending under a node still never exceeds that
node's allowance. The model's walk stops at the first node, counting from the
spender, that cannot take the amount, and a failed spend names it; that node
is where a refill request goes, up the delegation chain to the human who holds
the admin document. The proven statement about the name is weaker: it is a
node on the spender's chain that cannot take the amount (Theorem~\ref{thm:budget}).
Refilling and issuing a new root are admin operations; an agent can only
append a child node by delegating. Allowances are cumulative and never reset.

Figure~\ref{fig:budget} shows two runs from a root of 10 with two children of
8 each and nothing spent. In the first, after node~1 spends 8, node~2's spend
of 3 fails at the root, which has 2 left, and its spend of 2 succeeds. In the
second, node~1 delegates an allowance of 20, more than its own 8, to a new
node~3, whose spend of 9 is refused and whose spend of 8 succeeds, charging
node~3, node~1 and the root.

\begin{figure}[t]
\centering
\input{budget-tree}
\caption{Oversubscribed allowances. The \lean{example}s at the end of
\file{Budget.lean} check by evaluation the first refusal's error, and each
run's final ledger and log of successful spends; ``spent'' lists each node's
spent after the run, from node~0. A refused spend changes no
node.}\label{fig:budget}
\Description{A delegation tree: root with allowance 10, two children with
allowance 8 each, and a node with allowance 20 delegated under the first
child, with two runs of spends and their outcomes.}
\end{figure}

\section{From Rust to the Model}\label{sec:proofs}

The theorems are proven about the Lean translation of the evaluator the
boundaries run, not about a separate model alone. Figure~\ref{fig:pipeline}
shows the chain:
\begin{enumerate}
\item The evaluator, \file{crates/grant-eval}, is \code{no\_std} Rust with no
  dependencies and \code{forbid(unsafe\_code)}, written in a restricted subset.
\item Charon compiles the crate to LLBC and Aeneas translates the LLBC to Lean
  (\file{formal/AttenuatingGrants/Extracted/GrantEval.lean}, generated, never
  edited by hand).
\item Proofs in \file{formal/AttenuatingGrants/Equiv/} show that the
  translation computes the readable model on every input.
\item The same Rust compiles to \code{wasm32v1-none} for the grant signer and
  the host's component, and natively for the relay.
\end{enumerate}
The readable model is close to a reference implementation, not an independent
specification: the equivalence rules out faults of the Rust (indexing, loop
bounds, early returns, overflow, comparing values of different variants), but
whether a caveat's meaning is what an owner intends is a question about the
model, such as the positional semantics of \S\ref{sec:grants}.

\begin{figure*}[t]
\centering
\resizebox{\linewidth}{!}{\input{pipeline}}
\caption{The proof runs from the deployed Rust back to the readable model.
Solid green edges are proven in Lean; dashed orange edges are trusted.
Theorem~\ref{thm:total}'s step bound is about a counted copy of the model, not
the translation.}
\label{fig:pipeline}
\Description{Rust source flows through Charon and Aeneas to a Lean
translation, which equivalence proofs relate to the readable model where the
four theorems are proven; the same Rust compiles through rustc to wasm32v1-none
for the grant signer and the agent host's component, and natively for the
relay.}
\end{figure*}

\subsection{The restricted subset}

The evaluator is written so that every item translates with no opaque part.
Data is plain structs and enums. One trait, \code{Check<A>}, is used with
static dispatch only; there are no trait objects, closures, iterator adaptors
or derived equality, so enums compare through explicit \code{match} and byte
strings through one loop. Each loop is a single \code{while i < s.len()} over
one slice with an early \code{return}, the loop form Aeneas supports. Nothing
on the evaluation path can panic: slices are indexed only under
\code{i < s.len()}, argument indices that come from caveats go through
\code{get}, and budget amounts use \code{checked\_add}. The generic check is:
\begin{lstlisting}[language=Rust]
pub fn allows<A, C: Check<A>>(grant: &[C], action: &A) -> bool {
    let mut i = 0;
    while i < grant.len() {
        if !grant[i].holds(action) { return false; }
        i += 1;
    }
    true
}
\end{lstlisting}
The generated Lean builds against Aeneas's library with no \code{axiom} or
\code{opaque} declaration. CI re-runs Charon and Aeneas, pinned to one Aeneas
commit and the Charon and Rust toolchain its Nix flake fixes, and fails if the
result differs from the committed translation\claim{extraction-fresh}. Aeneas
runs with \code{-warnings-as-errors}, yet logs that ``the crate contains
extracted external, unknown definitions''. Reading Aeneas's source at the
pinned commit, that line is a log message, not one of the diagnostics the
option governs, and it is printed whenever the translation refers to a
definition with no Rust body, as it does for the library items Aeneas's Lean
library models (\code{Vec::push}, slice \code{get}, \code{usize::BITS}).

\subsection{Equivalence}

For each vocabulary $V$, \file{Equiv/V.lean} defines a denotation from the
translated types to the model's. It is part of the specification, so we show
chain's, verbatim:
\begin{lstlisting}[language=Lean]
def denoteVal : grant_eval.chain.Val → Val
  | .I128 x => .i128 x.bv
  | .U64 x => .u64 (denoteU64 x)
  | .Bool b => .bool b
  | .Bytes v => .bytes (bytesOf v.val)
  | .Address v => .address (bytesOf v.val)
  | .Symbol v => .symbol (bytesOf v.val)
  | .Other v => .other (bytesOf v.val)

def denoteAction : grant_eval.chain.Action → Action
  | .Invoke contract function args =>
    .invoke (bytesOf contract.val) (bytesOf function.val)
      (args.val.map denoteVal)
  | .CreateContract wasmHash =>
    .createContract (bytesOf wasmHash.val)

def denoteCaveat : grant_eval.chain.Caveat → Caveat
  | .ContractIn contracts =>
    .contractIn (contracts.val.map fun c => bytesOf c.val)
  | .FunctionIn functions =>
    .functionIn (functions.val.map fun f => bytesOf f.val)
  | .ArgI128AtMost index n => .argI128AtMost index.val n.bv
  | .ArgEq index val => .argEq index.val (denoteVal val)
\end{lstlisting}
and $\mathit{denoteGrant}(g)$ is the grant of $g$'s denoted caveats. Each case
maps a constructor to the constructor of the same name, and byte strings by a
proven-injective map (\lean{bytesOf\_inj}). The equivalence theorem for chain
reads, verbatim:
\begin{lstlisting}[language=Lean]
theorem extracted_allows (g : Slice grant_eval.chain.Caveat)
    (a : grant_eval.chain.Action) :
  grant_eval.chain.allows g a
    = ok ((denoteGrant g).allows (denoteAction a))
\end{lstlisting}
The same holds for files, network and signing, and for \code{allows\_all} over
a slice of actions\claim{equiv-chain}\claim{equiv-files}\claim{equiv-network}\claim{equiv-signing}.
Aeneas's \code{Result} is an interaction tree, and its \code{loop} combinator,
which every translated \code{while} calls, is defined in its library by
\code{partial\_fixpoint}. Termination of the translated functions is therefore
not a property of their definitions but a theorem: the right-hand side is
\code{ok}, which excludes both a failure (a panic, an overflow, an
out-of-bounds index) and divergence\claim{rust-total}. The proofs rest on two
loop lemmas, one for the ``any'' and one for the ``all'' shape, proven by
induction on the elements left, and on lemmas for byte-string comparison and
the path-prefix loop.

Two properties of the translation widen what the theorems cover. Aeneas gives
\code{usize} the width \code{System.Platform.numBits}, which proofs know only
to be 32 or 64, so the theorems hold for the 32-bit \code{wasm32v1-none}
builds although Charon ran on a 64-bit host. And Aeneas models arithmetic as
failing on overflow; the theorems show none occurs, and the workspace builds
with \code{overflow-checks = true} in its release and contract profiles, so
trapping and wrapping builds agree.

\subsection{The four theorems}

\begin{mainthm}[Totality and cost]\label{thm:total}
For every vocabulary, grant and action, the translated \code{allows} and
\code{allows\_all} return a result (above). For the step-counting Lean
evaluator of \S\ref{sec:cost}, with caveats $cs$, an action $a$ and a batch
$as$:
\begin{align*}
\mathit{steps}(cs, a) &\le 2\,\bigl(|cs| + \textstyle\sum_{c\in cs} |c|\bigr)\,(|a| + 1),\\
\mathit{steps}(cs, as) &\le 2\,\bigl(|cs| + \textstyle\sum_{c\in cs} |c| + 1\bigr)\,\textstyle\sum_{a\in as}(|a| + 1),
\end{align*}
and the counted evaluator returns the model's verdict.
\end{mainthm}
\leannames{\lean{V\.allows\_cost\_bound}, \lean{V\.allowsAll\_cost\_bound},
\lean{V\.allowsCounted\_fst} for each vocabulary
$V$\claim{cost-bound}\claim{cost-bound-batch}\claim{counted-agrees}.}
\noindent The bound is about a hand-written counted copy of the model, not the
translation: no theorem relates its steps to the extracted code, whose byte
comparison checks lengths first where the counted one walks pairwise. Its
product form is also loose. Each caveat's work is a membership scan bounded by
the caveat, an argument index bounded by the action, or a comparison bounded
by the shorter operand, so a sum of the form
$O(\sum|c| + |cs|\,(|a|+1))$ should bound it; that is stated, not proven. On
the measured points the bound exceeds the model's actual steps by up to a
factor of 33,000\claim{bound-looseness}, and priced at the measured 99.4
instructions per step it exceeds Soroban's per-transaction limit, so it says
the check terminates in polynomial time, not that it fits in a transaction.
\S\ref{sec:eval} measures what the deployed code does\claim{deployed-cost};
that its work is proportional to the step count on every input is not
proven\claim{deployed-cost-general}.

\begin{mainthm}[Attenuation]\label{thm:atten}
$(g.\mathit{attenuate}\ c).\mathit{narrower}\ g$, and likewise for any list of
added caveats.
\end{mainthm}
\leannames{\lean{attenuate\_narrows},
\lean{attenuateAll\_narrows}\claim{attenuation}.}

\begin{mainthm}[Chaining]\label{thm:chain}
Narrowness is transitive; every grant reached from a root by any number of
attenuation steps is narrower than the root; and a narrower grant accepts a
batch only if its ancestor does.
\end{mainthm}
\leannames{\lean{narrower\_trans}\claim{chaining},
\lean{Delegates\.narrower} and \lean{delegate\_narrows}\claim{delegation-chain},
\lean{allowsAll\_of\_narrower}\claim{batch-narrower}.}
\noindent In the model, Theorems~\ref{thm:atten} and~\ref{thm:chain} are
immediate consequences of conjunction, with one-line proofs. What the
equivalence adds is their carry-over to Rust caveat slices: a slice reached
from a root slice by pushing caveats accepts, alone or across a batch, only
what the root accepts (\lean{V\.extracted\_delegates\_narrower},
\lean{V\.extracted\_allowsAll\_of\_delegates})\claim{rust-delegation}. This is
the case both boundaries implement, since a child grant's stored caveats are
its parent's followed by the added ones.

\begin{mainthm}[Budget safety]\label{thm:budget}
For the translated \code{budget} module and every ledger $s$ with denotation
$L$: \code{spend} returns the model's outcome and ledger; a successful spend at
$i$ adds the amount to every node $a$ with $L.\mathit{Within}\ i\ a$ and
changes no other node; a failed spend returns $s$ unchanged; an
\code{Overdrawn a} error names a node on $i$'s chain that cannot take the
amount; and, for a Lean-defined caller that applies spends, refills and
appends in sequence from a ledger with nothing spent, the total successfully
spent at or under any node is at most that node's allowance at the end of the
sequence.
\end{mainthm}
\leannames{\lean{Budget\.extracted\_spend}\claim{budget-equiv},
\lean{extracted\_spend\_ok\_charges}\claim{budget-charges},
\lean{extracted\_spend\_ok\_unchanged}\claim{budget-unchanged},
\lean{extracted\_spend\_error}\claim{budget-failure},
\lean{extracted\_spend\_overdrawn}\claim{budget-overdrawn},
\lean{extracted\_spentUnder\_le\_allowance}\claim{budget-total},
all in \file{Equiv/Budget.lean}.}
\noindent The last statement, verbatim:
\begin{lstlisting}[language=Lean]
theorem extracted_spentUnder_le_allowance
    (v : alloc.vec.Vec budget.Node) (ops : List RustOp)
    (hcap : v.length + ops.length ≤ Usize.max)
    (h0 : ∀ n ∈ v.val, n.spent.val = 0) :
  ∃ v' log, runExtracted v ops = ok (v', log) ∧
    ∀ (a : Nat) (n : budget.Node), v'.val[a]? = some n →
      spentUnder (denoteLedger v'.val)
          (log.map fun e => (e.1.val, e.2.val)) a
        ≤ n.allowance.val
\end{lstlisting}
It has no hypothesis on the allowances, so it holds however much children
oversubscribe their parent; the allowance it bounds by is the final one, since
refills raise it. It covers runs that grow the tree: each logged spend is
measured against the final tree, and the proof shows an append changes no
existing parent link, so a spend lies under $a$ in the final tree exactly when
it charged $a$. \code{runExtracted} is a Lean definition that calls the
extracted \code{spend}, \code{refill} and \code{append} in sequence on one
ledger, as a Rust caller holding a \code{Vec} would; it is not extracted code.
The grant signer calls the budget functions differently, on a ledger of one
grant's lineage per authorization; \file{Signer.lean} follows those calls and
proves the same bound for every sequence of them (\S\ref{sec:soroban}).

Two side results discharge the Rust-specific risks. The charging loop's
unchecked \code{spent += amount} never overflows and visits exactly the chain
the check accepted (\lean{Budget\.charge\_loop\_eq})\claim{budget-no-overflow}:
the walk found $\mathit{spent} + \mathit{amount} \le \mathit{allowance} \le
2^{128}-1$ at each node, and each node is visited once. Amounts are unbounded
naturals in the model and \code{u128} in Rust; a refill fails past
$2^{128}-1$ as \code{checked\_add} does, and a spend whose sum overflows also
exceeds the allowance, so the two agree with no wrapped or saturated case.
\code{append} calls \code{Vec::push}, which Aeneas models as failing only when
a vector holds \code{usize::MAX} elements, so the append and run theorems
assume the ledger stays below that\claim{ledger-size}.

\subsection{The cost model}\label{sec:cost}

Each vocabulary has a counted evaluator that returns the model's Boolean paired
with a step count. A step is one primitive comparison (two bytes, two scalars,
two \code{i128}s by signed order), one constructor test, or one visit to a list
cell (a caveat, an action in a batch, a candidate in a membership list, an
argument while indexing, a path component); reaching the end of a list costs
nothing. Action and caveat sizes count the bytes and list cells a check can
touch: a chain invocation's size is its contract and function lengths plus
$\sum(|\mathit{arg}| + 1)$, a path's is $\sum(|\mathit{component}| + 1)$, and
\code{argI128AtMost} has size 2 whatever its index. The constant is 2 in every
vocabulary. The batch bound needs the $+1$ in its grant factor: visiting an
action's cell costs one step whatever the grant, so the empty grant checks a
batch of $n$ actions in exactly $n$ steps, which a factor of
$|cs| + \sum |c| = 0$ would deny. The model counts a call's fixed overhead
(entering the function, decoding the grant) as zero.

\subsection{What the build checks}

\code{lake build} checks every theorem above. CI rejects \code{partial},
\code{unsafe}, \code{sorry}, \code{admit}, \code{native\_decide},
\code{implemented\_by}, \code{extern} and \code{axiom} anywhere in the
package's Lean sources, the translation included\claim{no-unchecked-declarations};
Aeneas's library, which defines \code{loop} by \code{partial\_fixpoint} and
declares axioms of its own, is outside that check. \file{Axioms.lean} records
each theorem's \code{\#print axioms} output under \code{\#guard\_msgs}: every
theorem depends on at most \code{propext}, \code{Quot.sound} and
\code{Classical.choice}, and none on an axiom of Aeneas's
library\claim{axioms-builtin}. CI runs the Lean build and the extraction check
on every change to the evaluator or the proofs. Table~\ref{tab:loc} gives the
sizes.

\begin{table}[t]
\caption{Sizes of the evaluator and its proofs (non-blank lines; Rust also
without \code{//} comment lines), counted by the build.}\label{tab:loc}
\small
\begin{tabular}{@{}lr@{}}
\toprule
\textbf{Component} & \textbf{Lines} \\
\midrule
Evaluator, \file{crates/grant-eval/src} (Rust) & \LocGrantEval \\
Readable model: \file{Grant.lean}, \file{Vocab/} & \LocModel \\
Budget model and its theorems, \file{Budget.lean} & \LocBudget \\
Cost model and bounds, \file{Cost/} & \LocCost \\
Equivalence proofs, \file{Equiv/} & \LocEquiv \\
Generated translation, \file{Extracted/} & \LocExtracted \\
\bottomrule
\end{tabular}
\end{table}

The evaluator's own tests complement the proofs at the level the proofs
trust: they check agreement with a plain-Rust reading of the semantics and
attenuation on every grant of up to three caveats from a small universe per
vocabulary\claim{eval-exhaustive}, and every sequence of up to four spends and
refills on an oversubscribed forest against a reference\claim{budget-exhaustive}.
These run the compiled code that the proofs reach only through the trusted
translation.

\section{The Soroban Boundary}\label{sec:soroban}

On chain, an agent's grant lives in a \emph{grant signer}
(\file{contracts/grant-signer}): a contract deployed once per account, holding
for each grant a rule-eval rule set (a root grant's rules, the ones the
account's scoped document for its agents lowers to, and for a delegated
grant the attenuation steps it adds), its holder's ed25519 key, its parent,
its budget node and a revoked flag. The account is the Perch account Nido's
factory mints, unchanged.

\subsection{A delegated signer, not a policy}

A policy contract attached to an agent's context rule is the usual way to add a
check to an OpenZeppelin smart account~\cite{ozStellarAccounts}. It would be
the wrong place here, for two reasons. Perch accounts are document-only, and
their document schema has no member that names another policy contract. More
importantly, a policy is asked only about the contexts the submitter routed to
its rule; every other context's rule would have to be shown unsatisfiable by an
agent. The document schema does declare delegated signers, which the account
authenticates through CAP-0071's \code{delegate\_account\_auth}, and the host
then calls the delegate's \code{\_\_check\_auth} with the account's payload
and its whole context list\claim{cap71-contexts}. So the account's document
names the grant signer as a signer and makes it the only signer of the rules
agents act through:
\begin{lstlisting}[basicstyle=\ttfamily\scriptsize]
"signers": [
  {"id": "owner", "verifier": "<WebAuthn verifier>",
   "key": "<passkey>"},
  {"id": "grants", "address": "<grant signer>"}],
"rules": [
  {"name": "admin", "scope": {"type": "self-admin"},
   "principals": {"type": "all", "signers": ["owner"]}},
  {"name": "a-token",
   "scope": {"type": "contract", "address": "<token>"},
   "principals": {"type": "all", "signers": ["grants"]}}]
\end{lstlisting}
The grant signer sees every context of the authorization, including contexts
the submitter routed to other rules, so the guarantee does not depend on rule
selection, which a delegate's signature does not bind. The payload a holder
signs is the host's payload for the account's entry, which names the account,
so a holder's signature authorizes for that one account. Adding agents to an
account is a document change and one deployment; nothing in Perch or
OpenZeppelin changes.

\subsection{Checking an authorization}

Figure~\ref{fig:soroban} follows one agent authorization. The account's
\code{\_\_check\_auth} runs Perch's guards and then OpenZeppelin's rule
check, which authenticates the rule's signer, the grant signer, by
\code{delegate\_account\_auth}. The grant signer's \code{\_\_check\_auth}
receives the payload, a signature naming the grant and carrying the holder's
ed25519 signature, and every context, and refuses at the first failure:
\begin{enumerate}
\item load the signing grant and its ancestors, refusing if any is revoked;
\item verify the holder's signature over the payload;
\item refuse a contract creation; a context that calls the grant signer
  itself, unless it is the holder's own delegation; a context that calls the
  account; or a call on the budget asset other than \code{transfer} with a
  non-negative amount;
\item read every other context through rule-core's Soroban facade and
  require rule-eval's \code{authorize} over the grant's rule set, each
  context by the first rule that allows it, with the grant's cap records;
\item charge each budget-asset transfer to the grant and every ancestor with
  rule-eval's \code{spend}, refusing if any node would be overdrawn, and only then
  store the new \code{spent} values.
\end{enumerate}
The grant signer writes nothing before it refuses. A refusal fails the
account's authorization, and with it the call that asked for it; a calling
contract can catch that failure and go on, so the transaction need not fail.
Soroban then rolls back the failed call's storage writes and authorization
state, so either way the refusal changes nothing. The holder's ed25519 signature covers the host's payload,
not the grant id carried beside it, so a key that holds two grants produces
signatures usable under either; one key per grant avoids this.

\paragraph{What the grant signer sees.} The contexts it checks are those that
need the account's authorization. A call's effects that need none are outside
the tree and outside every caveat and budget: a \code{transfer\_from} against
an allowance the owner gave a contract earlier, a contract that already holds
or controls the account's funds, or any allowed call whose side effects never
call \code{require\_auth} for the account. The guarantee against hidden calls
holds for calls that need the account's authorization, and the assumption that
no allowed contract has such a standing position (\S\ref{sec:threat}) carries
the rest\claim{no-standing-allowance}.

\begin{figure}[t]
\centering
\resizebox{\linewidth}{!}{\input{soroban-flow}}
\caption{An agent's call through the account. The host passes every context of
the account's authorization to the grant signer, which checks them all with
the evaluator; one disallowed context refuses the whole tree.}
\label{fig:soroban}
\Description{An authorization tree in which an allowed router call nests a
token transfer; both contexts reach the grant signer, which runs rule-eval's authorize.}
\end{figure}

\paragraph{Delegation, and the grant signer's second role.}
\code{delegate(parent, key, steps, allowance)} stores a child grant whose
rule set is the parent's with the given rule-eval attenuation steps applied,
with a budget record under the parent's. It requires the grant signer's own authorization, so the
host calls the grant signer's \code{\_\_check\_auth} in two roles, with the
same three arguments: as the account's delegate, with the account's payload
and contexts, and as the check of the grant signer's own address, whenever a
call needs that address's authorization (its own \code{delegate}, a call on
another contract that asks for it, a deployment with it as deployer). The
host lets a contract's \code{require\_auth} of itself call its own
\code{\_\_check\_auth} and refuses any other re-entry. The checks are the
same in both roles: contract creation is refused under every
grant\claim{signer-deploy}, since in the second role no account document
stands behind the check; a call on the account, an unmetered call on the
budget asset and a call on the grant signer other than the holder's own
delegation are refused; and the rule set bounds every other call. A
signature in the second role authorizes nothing for the account, which
accepts the grant signer only through
\code{delegate\_account\_auth}\claim{signer-own-address}. One consequence is
stated rather than enforced. Any contract can call \code{require\_auth} on an
address it is given, so a holder can have a contract its grant allows ask
for the grant signer's own authorization, and authorize what that contract
does as the grant signer's address, within the grant's rules. That includes
moving what the address holds; a transfer of the budget asset is still
charged to the grant, and one of any other asset is bounded only by the
rules.
The deployment gives the address no balance: a contract's balance in a
Stellar Asset Contract, the native one included, is absent until something
is sent to it, and storage rent is paid in fees by the transaction's
submitter. But the address is published in the account's document, and
whatever is sent to it is reachable this way.

\subsection{Threat coverage}

Table~\ref{tab:soroban-threats} maps each attempt in \S\ref{sec:threat} to
what refuses it. Each refusal is asserted in an integration test by the
contract that refused and its error, on a factory-minted Perch account under
enforcing authorization, where the host runs every \code{\_\_check\_auth}.

\begin{table*}[t]
\caption{Escalation attempts at the Soroban boundary and what refuses them.}
\label{tab:soroban-threats}
\small
\begin{tabularx}{\linewidth}{@{}>{\raggedright\arraybackslash}p{0.27\linewidth}>{\raggedright\arraybackslash}X l@{}}
\toprule
\textbf{Attempt} & \textbf{What refuses it} & \textbf{Status} \\
\midrule
Issue, revoke, rotate or refill a grant, its own included &
  These need the account's authorization, which only the owner's admin rule
  gives; no agent rule is scoped to the grant signer, and asked directly the
  grant signer refuses every call on itself but the holder's own delegation. &
  \claim{signer-admin} \\
Change the account's rules or policies, or call its admin functions &
  No agent rule is scoped to the account; a document that wrongly scopes one
  still meets the grant signer's refusal of calls on the account. &
  \claim{signer-account} \\
Delegate from an ancestor's grant, or one it does not hold &
  The grant signer accepts a delegation only with the signing grant as parent. &
  \claim{signer-delegate-other} \\
Widen a grant by delegating &
  A child's rule set is its parent's with rule-eval attenuation steps
  applied, each of which narrows. &
  \claim{signer-narrows}\ \claim{signer-rule-eval-chain} \\
Hide a disallowed call inside an allowed one &
  Every context of the account's authorization is checked; a nested call the
  signed tree leaves out fails its \code{require\_auth} in the host. Effects
  that need no authorization from the account are outside the tree. &
  \claim{signer-hidden-call}\ \claim{batch-narrower} \\
Deploy a contract &
  The grant signer refuses every contract creation, under every grant and in
  both roles; through the account, a Perch document has no creation scope,
  so the account refuses it first. &
  \claim{signer-deploy}\ \claim{create-contract-denied} \\
Hand off the budget asset (\code{approve}, \code{burn}) &
  Only metered transfers of the budget asset pass. &
  \claim{signer-metered} \\
Overspend through siblings or several transfers &
  Every transfer in the authorization is charged to the spender and every
  ancestor; the first overdrawn node refuses and nothing is stored. &
  \claim{signer-oversubscribed}\ \claim{signer-all-charged}\ \claim{signer-budget-total} \\
Keep acting after revocation, or with a rotated-out key &
  Revocation of any ancestor refuses; rotation replaces the key. &
  \claim{signer-revocation}\ \claim{signer-rotation} \\
\bottomrule
\end{tabularx}
\end{table*}

\subsection{The wiring}

The evaluator is verified; what connects it to Soroban is not. The grant
signer reads each context through rule-core's Soroban facade, as the relay
reads a call's XDR: a contract and an address argument by its
\code{ScAddress} XDR, a function and a symbol by its bytes, \code{i128},
\code{u64}, \code{bool}, bytes and strings by value, and any other argument
as \code{other} with its full XDR. A test checks that the grant signer and
the relay read every argument type alike\claim{signer-relay-mapping}; that
the reading keeps distinct arguments distinct rests on XDR encoding being
injective and is not proven\claim{xdr-injective}. On every grant that names
its contracts, the rule set a caveat list lowers to allows exactly what the
caveats allow\claim{chain-rule-eval-agree}, and the grant signer reproduces a
recorded corpus of the outcomes it had when it held caveat lists and checked
them with the chain evaluator\claim{signer-check-auth-corpus}.

\paragraph{The budget link.} The grant signer never holds the ledger of
every grant whole. Its storage holds one grant per id, whose parent,
allowance and \code{spent} are that grant's node; \code{issue} and
\code{delegate} store a grant at the next id with nothing spent, which is the
model's \code{append}; \code{refill} calls rule-eval's \code{refill} on a
ledger of the one grant; and each \code{\_\_check\_auth} builds a ledger of only
the signing grant's lineage, root first with each node's parent the one before
it, calls rule-eval's \code{spend} on it once per budget-asset transfer, stopping
at the first failure, and only if every spend succeeded writes the changed
\code{spent} values back. \file{formal/AttenuatingGrants/Signer.lean} follows
these steps. Spending on the lineage ledger has the whole ledger's outcome at
the signing grant, an error naming the same grant, and afterwards the whole
ledger's lineage, and the spend loop over grant-eval's extracted \code{spend}
computes it\claim{lineage-ledger}. The contract's rule-eval \code{spend}
agrees with that extracted \code{spend} on lineage ledgers, on every outcome,
overdrawn grant and ledger, for amounts up to the largest \code{i128}, which
\code{issue}, \code{delegate} and \code{refill} keep every allowance
within\claim{signer-rule-eval-budget}. One authorization's spends are all or nothing:
the contract stores them only after every spend succeeds, and it does so in
the \code{require\_auth} of the authorization's root, inside the call that
encloses every call the authorization allows. When that call fails, whether
or not a caller catches the failure, the host rolls back its storage writes,
the charge included, and a charge stays only with the root call it was made
in. A test catches a refused authorization and an authorization whose
nested transfer fails after the charge, and finds the budget unchanged both
times\claim{host-frame-rollback}. So every sequence of the grant signer's calls is a run of the model on the
ledger of every grant\claim{budget-batch-as-run}, and Theorem~\ref{thm:budget}'s
bound holds for it: the total spent at or under any grant never exceeds that
grant's allowance, with no hypothesis on the
allowances\claim{signer-budget-total}. What remains trusted is that the
contract's few lines of storage, lineage and loops are those
definitions\claim{wiring-budget-calls}. The tests check the outcome:
oversubscribed siblings cannot together pass their parent, and a refused
spend stores nothing\claim{signer-oversubscribed}.

Two requirements on the account's document are stated and met by the
documents the tests apply, but nothing on chain checks them: a holder key must
reach the account only through the grant signer, since registering it as a
signer by another route lets it authorize without its grant; and the document
must name its own account's grant signer, since the grant signer cannot see
which account delegated to it\claim{document-requirements}.

\section{The Agent Boundary}\label{sec:host}

An agent's tool calls reach the file system, the network and the key keeper
through the agent host (\file{crates/agent-host}), which checks each call
against the agent's grant before acting on it. The host is an in-process
library, not a sandbox for the agent: nothing in it confines the agent's
process, which is why its being the agent's only path to those resources is
an assumption about the deployment\claim{sole-path}, and its \code{call}
takes the caller's agent id as given, so a deployment must bind each agent's
tool-call channel to its id\claim{agent-identity}. Process sandboxes such as
nono~\cite{nono} or a verified runtime such as WaVe~\cite{wave} are what
confine an agent; the host decides what a confined agent's requests may do.

The host has two interfaces. The admin interface issues a root agent with a
grant and revokes an agent together with every agent under
it\claim{host-revocation}. The agent interface takes five tool calls:
\code{read\_file}, \code{write\_file}, \code{fetch}, \code{sign} and
\code{delegate}. None reads or edits a grant\claim{host-no-grant-tool};
\code{delegate} starts a sub-agent whose grant is the caller's caveats
followed by the added ones, so every child allows a subset of its
parent\claim{host-delegation}.

\subsection{The evaluator as a component}

The verdict comes from \file{crates/grant-component}: \code{grant-eval}
compiled for \code{wasm32v1-none}, the target the grant signer uses, and
wrapped as a WebAssembly component behind a WIT interface with one interface
per vocabulary, each exporting \code{allows} and \code{allows-all}. The
component imports nothing and keeps no state between calls, so it can compute
only on the grant and action the host passes in. The host builds the component
at compile time, instantiates it in Wasmtime with an empty linker, and passes
the grant and the action on every call. Since WIT has no 128-bit integer, a
chain \code{i128} crosses the boundary as a pair of 64-bit halves. Tests check
that the component returns the native evaluator's verdict on every grant of up
to three caveats from a small universe per vocabulary, and for chain on every
authorization tree of up to two contexts\claim{component-agrees}; that
Wasmtime and the bindings carry values across faithfully is
trusted\claim{runtime-trusted}.

The evaluator is proven total and pure, so the component does not defend
against the evaluator's logic. It does two other things. It runs the build the
chain runs: the same source for the same \code{wasm32v1-none} target, so the
host's verdicts come from the artifact class the grant signer deploys, not
from a separate native build. And it confines that build, which is the output
of an unverified toolchain, to computing a verdict: a component that imports
nothing can touch no file, socket or key whatever its compiled code does. The
price is the boundary crossing measured in \S\ref{sec:eval}.

\subsection{Building actions}

The wiring that turns a tool call into an action is the trusted part of the
host, kept to three modules.

\paragraph{Files.} A requested path must be absolute. Each existing prefix is
canonicalized with \code{realpath}, which follows symlinks, applies \code{..}
to the directory a symlink led to, and returns the names stored on disk;
components past the last existing one are appended as given. Grant prefixes
are resolved the same way when a grant is issued or delegated. On a
case-insensitive file system the check therefore sees the stored name, and a
differently cased path is allowed exactly when the stored path
is\claim{host-case}. The host checks the canonical path, opens the file with
\code{O\_NOFOLLOW} on its last component, and checks the path of the opened
descriptor again (\code{F\_GETPATH} on macOS, \code{/proc/self/fd} on Linux)
before any data moves through that descriptor, so a component swapped for a
symlink between check and open is caught\claim{host-swap}. A new file is
created under its directory's checked descriptor with
\code{O\_CREAT|O\_EXCL|O\_NOFOLLOW}. Neither \code{..} nor a symlink reaches
outside the grant\claim{host-dotdot}\claim{host-symlink}. Three cases are not
covered. A file renamed or relinked by another process after the descriptor's
path was checked is read or written all the same, since the descriptor follows
the inode, not the path. A hard link inside the grant to a file outside it is
allowed, since the vocabulary names paths, not files. And an error while
resolving is reported before the check, so it can reveal that something
exists at a path the grant does not cover. The agent's own tool calls create
neither renames nor links.

\paragraph{Network.} The action is the host name the agent requested, as the
URL parser normalizes it (lower-case ASCII, IDNA names in \code{xn--} form), with
its port and scheme. After the check, the host's client resolves the name and
connects directly; the agent never supplies an address, and what a name
resolves to is trusted\claim{os-wiring}. A denied host, port or scheme is
refused before any connection\claim{host-network-denied}. The client follows
no redirects itself: each \code{Location} is resolved against the current URL
and checked as a new action before the next request, up to five hops, so a
redirect to a disallowed host is refused without contacting
it\claim{host-redirect}. The vocabulary binds names, not addresses: an
allowed name that resolves, or is rebound, to a loopback, link-local or cloud
metadata address is reached, and the HTTP method and path are not part of the
action. A caveat on the resolved address's class, or refusing private ranges
unless named, would close most of the first.

\paragraph{Signing.} The agent names a key id, a network passphrase and a
payload, and holds no key. A grant carries a chain caveat list beside its
signing list, parsed and mapped by the relay's code. The host checks the
signing action (key, network, payload kind) first; for a transaction or an
authorization entry it then decodes the payload and checks every context the
signature authorizes against the chain grant with \code{allows\_all}, and
only then calls the key keeper, so a denied request never reaches
it\claim{host-signing-before-keeper}. For a transaction the contexts are its
host function and every context of every authorization entry it carries; for
an authorization preimage, the account's own or a delegated signer's, they
are every context of its invocation tree\claim{host-signing-chain}. The host
signs only what the chain vocabulary describes: an envelope with exactly one
\code{InvokeHostFunction} operation, and no fee bump, second operation,
other operation or contract-code upload\claim{host-signing-refuses}. The
keeper is given the encoding of the decoded value the host checked, so the
bytes signed are the bytes checked\claim{host-signs-checked}. Each payload
kind is also bound to its network: a transaction is signed as its hash on the
named network; a preimage is signed only after the host checks that its
network id is the hash of the named passphrase\claim{host-delegated-preimage};
and a message is signed under SEP-53's prefix ``\code{Stellar Signed
Message:\textbackslash n}''~\cite{sep53}, which keeps a message signature from
being a transaction or authorization-entry signature. A message signature is
not bound to any network, so for messages a \code{networkIn} caveat restricts
only the label the agent supplies.

No caveat covers a transaction's fee, sequence number, time bounds, memo or
source account, an authorization entry's nonce or expiration ledger, or which
account a key id stands for, so an agent allowed to sign transactions can have
its key pay any fee for an allowed call. An empty chain list allows every
contract call, as an empty grant does, so a grant that allows
\code{transaction} or \code{authEntry} signatures should carry chain caveats.

\section{The Relay and Custody Signing}\label{sec:relay}

Agents that hold no keys submit Soroban calls through a relay
(\file{crates/grant-relay}) that speaks the protocol Nido's browser client
already uses: the client posts a host function and its signed authorization
entries to \code{/relay} and polls for the result, so a deployment changes only
the client's relay URL\claim{relay-frontend}. The relay runs the native build
of the same evaluator. Its grant bounds what the custody key will sign, pay
for and submit; authority over an account is still decided on chain, by the
account's own \code{\_\_check\_auth} and, for agents, its grant signer. The
repository's other relay, \file{infra/relayer}, wraps a third-party
fee-sponsoring plugin with no grant check and is outside this design: it
relays any contract call whose on-chain authorization is valid, which for an
account with a grant signer still means the grant signer's check. Each agent has a chain grant over the contexts of the
call's authorization trees and a signing grant over the custody signature.
Before any key is used, the relay requires, in order: that the API key names
an agent; that no supplied entry uses the transaction source's credentials;
that the chain grant, as the rule-eval rule set its caveats lower to, allows
the host function and every context of every supplied tree, read through
rule-core's Soroban facade as the grant signer reads them; and that the
signing grant allows the custody key, the
network and the \code{transaction} payload kind. These checks need no network,
so a denial leases nothing and never reaches the key
keeper\claim{relay-nested}.

It then leases a channel account and runs two simulations. A recording
simulation of the bare host function must succeed, must not ask for the
channel's own authorization\claim{relay-channel-auth}, and every context it
records must pass the chain grant\claim{relay-recorded}; an enforcing
simulation with the supplied entries must succeed, and the fee, the simulated
minimum plus a margin for rent that grows before the transaction lands, must
stay under a cap\claim{relay-fee-margin}. The supplied trees are what the
user's signatures authorize, and the recorded trees are what the execution asks for, so a call is relayed only when
the grant allows every context in either. That the two together cover every
context the signatures authorize rests on the host's matching of
\code{require\_auth} to signed trees and is not proven\claim{trees-cover}. The
chain grant does not see calls that ask for no authorization, calls a contract
makes inside its own \code{\_\_check\_auth}, or the fee, which the cap bounds.
Soroban RPC's recording simulation records authorization only for the root
invocation, so the relay cannot submit a call whose first request for an
account's authorization is made in a nested call\claim{network-root-recording}.

\paragraph{Relay grants and chain grants.} The relay's grants are a
configuration file with no link to the grant signer, so a deployment must
rewrite it when the admin document issues, revokes or rotates grants on chain;
that a deployment keeps the two in step is unproven\claim{relay-grants-in-step}.
A relay grant wider than the chain's lets more calls reach the network's
simulation and the grant signer, which still refuse them: with relay grants
allowing every call, all 44 attempts the network runs are
refused\claim{network-attack-suite}, each but one by the component that refuses
it in the simulation\claim{network-as-simulation}. A relay grant narrower than
the chain's refuses calls the account would allow.

\paragraph{Channels and the key keeper.} Transactions go out from a pool of
channel accounts, each with master weight 0 and the custody key as its only
signer; the relay checks this for every channel at start. A request holds its
channel until its transaction settles, so concurrent requests use distinct
channels, the relay reads each channel's sequence number fresh, and each
channel carries at most one transaction per ledger\claim{relay-throughput}.
The relay has no budget, so concurrent spends that together overdraw an
ancestor all pass it; the grant signer refuses the excess on chain, and the
refused transactions cost their channels a fee\claim{network-race}. The key
keeper signs the transaction hash with the custody key and the relay verifies
the signature before attaching it. The keeper is an interface with three
implementations: an in-process ed25519 keeper, a client for the Dfns Keys
API's Generate Signature endpoint with Dfns's user-action signing flow, and a
local mock of that endpoint that follows Dfns's OpenAPI document. Dfns's own
policies remain a check behind the grant, but they see only a 32-byte hash,
which their amount and recipient rules cannot evaluate~\cite{dfnsSigningPolicies}.

\paragraph{Channel accounts on testnet.} A channel's source account is not the
custody key's own address, so the arrangement depends on the network accepting
a transaction whose only signature is a signer's that is not the source. On
Stellar testnet, with the custody key held by the Dfns mock, channel-sourced
classic payments, Soroban invocations and fee bumps signed only by the custody
key were accepted, and the same transactions signed by a channel's zeroed
master key or by an unrelated key were refused with
\code{txBAD\_AUTH}\claim{testnet-channel}. Whether real Dfns signs such
transactions is open: its reference does not say whether the
\code{Transaction} signature kind checks the envelope's source account against
the key's own address\claim{dfns-transaction-source}, or whether a
\code{Hash} request on an ed25519 key signs the 32 bytes directly as the
message\claim{dfns-hash}.

\section{What Is Proven, Tested and Trusted}\label{sec:tcb}

Table~\ref{tab:tcb} separates the parts of the system by how each is assured.
The proofs cover the evaluator's translation, from \code{allows},
\code{allows\_all} and the budget functions to the model; everything that
feeds the evaluator actions or acts on its verdict is wiring, tested but not
proven.

\begin{table*}[t]
\caption{Assurance by component. Lines are non-blank Rust lines that do not
start with \code{//}, counted by the build.}\label{tab:tcb}
\small
\begin{tabularx}{\linewidth}{@{}>{\raggedright\arraybackslash}p{0.21\linewidth}r>{\raggedright\arraybackslash}X@{}}
\toprule
\textbf{Component} & \textbf{Lines} & \textbf{Assurance} \\
\midrule
Evaluator, \file{crates/grant-eval} & \LocGrantEval &
  Translation proven equal to the model on every input; Theorems 1--4
  (\S\ref{sec:proofs}). Exhaustive tests on small universes. \\
Grant signer, \file{contracts/grant-signer/src} & \LocSigner &
  Wiring: context mapping (\file{actions.rs}, \LocSignerActions), lineage,
  guards and storage (\file{contract.rs}, \LocSignerContract). Integration
  tests and the attack suite under enforcing authorization. \\
Component, \file{crates/grant-component} & \LocComponent &
  Argument conversion only; agreement with the native evaluator tested. \\
Agent host, \file{crates/agent-host/src} & \LocHost &
  Wiring for files (\file{fs.rs}, \LocHostFs), network (\file{net.rs},
  \LocHostNet) and signing (\file{sign.rs}, \LocHostSign), plus grants, the
  host loop, the workload runner and the benchmark. Tested, including symlink,
  case and check-then-open races, and by the host attack suite. \\
Relay, \file{crates/grant-relay/src} & \LocRelay &
  XDR-to-action mapping (\file{actions.rs}, \LocRelayActions), configuration,
  the request pipeline, channels and RPC. Tested in process, against Nido's
  TypeScript client, and on a local Stellar network. \\
\midrule
\multicolumn{3}{@{}l}{\textbf{Trusted, not proven}} \\
\multicolumn{3}{@{}>{\raggedright\arraybackslash}p{\linewidth}@{}}{%
  rustc for the host and \code{wasm32v1-none}; Charon and Aeneas; Aeneas's
  Lean library, including its interaction-tree \code{Result}, its \code{loop}
  defined by \code{partial\_fixpoint}, and its models of the Rust items the
  translation calls; Lean's kernel\claim{toolchain-trusted}. The denotation
  functions, which are part of the specification (\S\ref{sec:proofs}).
  Wasmtime, \code{wit-bindgen} and \code{wit-component}\claim{runtime-trusted}.
  The Soroban SDK and host, whose CAP-0071 behaviour Perch's tests pin for the
  pinned version\claim{cap71-contexts}. The operating system's
  \code{realpath}, descriptor paths and DNS\claim{os-wiring}. The key
  keeper\claim{keeper-trusted}. The Soroban RPC server's simulations and
  status reports, on which the relay's recorded-tree check and fee depend.} \\
\bottomrule
\end{tabularx}
\end{table*}

The trust in the translation is narrower than it looks: it covers Aeneas's
models of the Rust \code{core} and \code{alloc} items the evaluator calls,
since the evaluator calls nothing else: \code{Vec} with \code{deref},
\code{deref\_mut}, \code{len} and \code{push}; slices with indexing,
\code{get}, \code{get\_mut} and \code{len}; checked \code{usize} and
\code{u128} addition and unchecked \code{u128} addition; \code{u32 as usize};
scalar comparison; and \code{usize::BITS}, read by a compile-time assertion
the proofs do not use. The same evaluator source reaches all three
deployments, but each deploys its own compiled artifact: the grant signer
embeds it in a Soroban contract built by rustc for \code{wasm32v1-none}, the
agent host runs a separately built \code{wasm32v1-none} component, and the
relay links the native build.

\paragraph{Known gaps in the wiring.} Table~\ref{tab:gaps} collects the gaps
the paper states, each in the section that describes it.

\begin{table}[t]
\caption{Known gaps in the current design.}\label{tab:gaps}
\small
\begin{tabularx}{\linewidth}{@{}>{\raggedright\arraybackslash}X>{\raggedright\arraybackslash}p{0.12\linewidth}@{}}
\toprule
\textbf{Gap} & \textbf{Where} \\
\midrule
Under an empty grant, the grant signer's own address can authorize any call
on another contract. & \S\ref{sec:soroban} \\
The holder's signature does not cover the grant id. & \S\ref{sec:soroban} \\
That the grant signer's storage, lineage and loops are the Lean definitions
of its budget path is trusted. & \S\ref{sec:soroban} \\
Effects that need no authorization from the account are outside the tree. &
\S\ref{sec:soroban} \\
No caveat binds a signed transaction's fee, sequence number, time bounds or
memo. & \S\ref{sec:host} \\
Network caveats bind names, not resolved addresses. & \S\ref{sec:host} \\
Renames after the descriptor check, and hard links. & \S\ref{sec:host} \\
\bottomrule
\end{tabularx}
\end{table}

\section{Evaluation}\label{sec:eval}

We ask four questions. Do the boundaries refuse the escalations of
\S\ref{sec:threat} without changing state? Does the deployed evaluator's cost
follow the proven step count? What does enforcement cost at each boundary?
Does custody signing over channel accounts work on a public network?

\paragraph{Method.} Every number below comes from a recorded run, and each
claim tag names the result file it is read from; this paper ran no new
experiments. The Soroban runs execute every contract from its wasm in
\code{soroban-env-host} 27.0.1 under enforcing authorization: Perch's account,
\code{WebAuthn} verifier, document compiler and interpreter as deployed on
testnet, Nido's factory and grant signer, and a test dApp. Each agent call is
one \code{AddressWithDelegates} entry carrying the holder's signature, so every
account's \code{\_\_check\_auth} runs. Soroban CPU instructions and memory are
the host's metered costs, the same on every machine and every run, and they
are the costs of a network running that host version; no network was measured.
Agent-host times are wall clock, from one recorded run on an Apple M5 Max
under macOS 26.5.2 with the build profile, scheduling class, power source and
load recorded; Table~\ref{tab:host-cost} gives the method. The testnet result
is a recorded run of a program that creates its own accounts and keys.

\subsection{Refusals}

The attack suite gives every attempt to the agent the owner authorized, and
snapshots state before and after it: at the Soroban boundary the balances of
the account, the recipients, the attacker and the dApp instances in both
assets, every grant (holder key, caveats, allowance, spent, revoked), and the
account's revision, document hash and pending upgrade; at the host the file
tree's hashes and symlink targets, a second server's request count, the key
keeper's signature log and every agent's verdict on two probe files. A refusal
passes only if the snapshots are equal. Controls are allowed calls of the same
shape as the attack next to them.

\begin{table}[t]
\caption{Attack suite: attempts, refusals with state unchanged, and allowed
controls.}\label{tab:attacks}
\small
\setlength\tabcolsep{4pt}
\begin{tabular}{@{}llrrr@{}}
\toprule
\textbf{Boundary} & \textbf{Category} & \textbf{Attempts} & \textbf{Refused} & \textbf{Controls} \\
\midrule
Soroban & self-escalation & 16 & 16 & 0 \\
 & nested-call smuggling & 8 & 8 & 1 \\
 & delegation workarounds & 18 & 18 & 6 \\
 & oversubscription & 3 & 3 & 2 \\
Host & self-escalation & 10 & 10 & 0 \\
 & delegation workarounds & 21 & 21 & 3 \\
\bottomrule
\end{tabular}
\end{table}

In Soroban simulation all 45 attempts are refused with state unchanged and
the 9 controls pass\claim{attack-suite-soroban}; at the host all 31 are, and
its 3 controls pass\claim{attack-suite-host} (Table~\ref{tab:attacks}). The
Soroban attempts include the grant signer's admin functions through the
account and as the grant signer itself; the account's \code{apply\_doc},
\code{execute} and \code{schedule\_upgrade}; documents that wrongly scope an
agent rule to the account or to the grant signer; transfers to the attacker
nested two to five calls deep behind an allowed call, or left out of the
signed tree; contract deployment; \code{approve} and \code{burn} of the
budget asset; arguments of another type with the same bytes (a muxed
recipient, a \code{u64} amount, a \code{u32} for the bytes of its XDR, a
string for a symbol); revoked and rotated-out keys; a replayed entry; and
siblings and a grandchild spending past an ancestor's allowance. Refusals come
from three places: the account (no rule covers the call), the grant signer
(\code{NotHolder}, \code{AccountCall}, \code{ContractCreation},
\code{UnmeteredCall}, \code{NotAllowed}, \code{Overdrawn}, \code{Revoked}, a
failed holder signature), and the host (no entry covers the call, a used nonce, a
re-entry). The host attempts are tool calls that issue or edit grants,
sub-agents using a wider grant than their parent's, revoked agents, paths that
name the work directory and reach outside it, URLs that name the allowed host
and reach another (another port or scheme, an address literal, a trailing
dot, userinfo, redirects), and signing requests for another key, network or
payload kind. These are simulation and host results; the attack suite has not
been run through a local Stellar network or through the relay.

The suite found two things. A muxed address (\code{M\ldots}) of a recipient
pinned by \code{argEq} is refused although it names the same account's key,
since values compare exactly. And it found a defect in the host's wiring: the
host signed only the account's own authorization preimage, not the
delegated-signer preimage a grant-signer holder signs, so an agent whose holder
key the keeper held could not authorize through its grant signer. The host
now signs both after the same network check\claim{host-delegated-preimage}. The
defect failed closed, and it sat in the unverified wiring, where
Table~\ref{tab:tcb} puts the remaining risk.

\paragraph{Workloads.} Two scripted agent sessions with sub-agents run in the
simulation, each call one transaction through the account: \emph{payments},
with a per-call cap of 100 and an allowance of 400, and a payroll sub-agent
pinned to one payee; and \emph{defi-router}, whose calls nest the account's
transfers up to four contexts deep. Every call has the expected outcome, and
every refused one leaves state unchanged. The router session shows the
positional limit of \S\ref{sec:grants}: its calls carry amounts at different
argument positions, so it bounds amounts by its allowance alone.

\subsection{The cost-bound cross-check}

Theorem~\ref{thm:total}'s bound is about a counted model. To test whether the
deployed evaluator follows the model's step count, the harness measures the
grant signer's \code{\_\_check\_auth} on five sweeps, each varying one thing:
the number of caveats; the length of a byte-string argument, compared by an
\code{argEq} and not compared; the number of contexts; and the length of a
\code{contractIn} list. Every point is measured twice on the same stored
grant: passing, with every caveat holding, and failing, with the first caveat
failing on the first context, so the evaluator stops at once. Loading the
grant, verifying the signature and mapping contexts to actions cost the same
in both, so the difference is the evaluator's own work. For every point the
harness writes the grant and actions in the evaluator's form, and the Lean
package evaluates its own counted evaluator and bound on them.

The contract, the native evaluator and the Lean model agree on every verdict.
Over the 28 pairs, the evaluator's work is 96 to 164 CPU instructions per Lean
step, and a least-squares fit through the origin of the instruction
difference on the step difference gives 99.4 instructions per step with
$R^2 = 0.986$\claim{deployed-cost}. The fit is through the origin over
differences, so it prices evaluator work only; the mapping of contexts to
actions, which grows with argument size, is outside it: with an argument that
no caveat compares, the whole check grows from 1.072M to 1.255M instructions
between 0 and 16 KiB while the model's steps stay at 68. Priced at that
constant, the bound covers the whole check above the zero-caveat cost on every
point, by a factor of more than 12, and the margin widens as arguments,
contexts and lists grow, since the bound grows with the product of grant and
action sizes and the measured cost with their sum\claim{bound-looseness}. That
the evaluator's work is proportional to the step count on every input remains
unproven\claim{deployed-cost-general}.

\subsection{Cost on chain}

Table~\ref{tab:signer-cost} gives the grant signer's whole
\code{\_\_check\_auth} on one context, from the caveats sweep: instantiating
the contract, loading the grant, checking the signature, mapping and
evaluating. The zero-caveat check costs 972,392 instructions, which the cost
model counts as zero. Under 4 caveats each further context costs about 97,000
instructions\claim{soroban-check-cost}. Whole agent transactions through the
account cost 2.17M to 2.69M instructions with one context and 5.04M with four
nested contexts, against a limit of 400M per
transaction\claim{soroban-tx-cost}.

\begin{table}[t]
\caption{Grant signer \code{\_\_check\_auth} on one context, every caveat
holding (CPU instructions and memory bytes).}\label{tab:signer-cost}
\small
\begin{tabular}{@{}rrrr@{}}
\toprule
\textbf{Caveats} & \textbf{CPU} & \textbf{Memory} & \textbf{Lean steps} \\
\midrule
0 & 972{,}392 & 416{,}582 & 1 \\
1 & 1{,}024{,}459 & 420{,}766 & 60 \\
4 & 1{,}168{,}746 & 434{,}627 & 188 \\
16 & 1{,}735{,}302 & 489{,}255 & 700 \\
64 & 4{,}003{,}738 & 706{,}951 & 2{,}748 \\
\bottomrule
\end{tabular}
\end{table}

\subsection{Cost in the host}

Table~\ref{tab:host-cost} times one check three ways: native
\code{grant-eval} on values built once; the host-side check without
WebAssembly, converting the host's WIT values natively and evaluating; and
the component through Wasmtime. Each case is a three-caveat grant and an
action it allows; the chain case checks a three-context tree. The component
check takes 0.30 to 1.01\,\textmu s, 15 to 47 times native and about 3 times
the host-side check without WebAssembly, so the price of running the deployed
build in a component is about a factor of 3 over checking natively in the
host; compiling and instantiating the component takes 22.0\,ms, once per
host\claim{host-check-time}. A whole tool call costs more than its check: a
denied fetch 0.68\,\textmu s, a read resolved and checked 37 to
47\,\textmu s, a signature about 10\,\textmu s, and an allowed fetch from a
local server 160\,\textmu s\claim{host-tool-call-time}. The load average was
about 5.5 during the run, and macOS cannot pin a thread to a core, so these
are the medians of one run, not a bound.

\begin{table}[t]
\caption{Per-check time (ns): medians with 95\% confidence intervals over 101
interleaved rounds, after a 200\,ms warm-up per operation, timed in batches of
at least 2\,ms; ratios within a round.}\label{tab:host-cost}
\small
\setlength\tabcolsep{4pt}
\begin{tabular}{@{}lrrrr@{}}
\toprule
\textbf{Case} & \textbf{Native} & \textbf{No Wasm} & \textbf{Component} & \textbf{Ratio} \\
\midrule
files & 12.5 & 182.5 & 579.9 & 46.5$\times$ \\
 & {\scriptsize 12.4--12.6} & {\scriptsize 180.7--183.6} & {\scriptsize 579.1--581.2} & \\
network & 7.6 & 90.4 & 298.6 & 39.6$\times$ \\
 & {\scriptsize 7.5--7.6} & {\scriptsize 89.9--90.7} & {\scriptsize 298.0--299.2} & \\
signing & 15.9 & 127.7 & 391.5 & 24.6$\times$ \\
 & {\scriptsize 15.8--16.0} & {\scriptsize 126.8--128.3} & {\scriptsize 390.8--392.3} & \\
chain (3 ctx) & 67.9 & 341.6 & 1{,}006.7 & 14.9$\times$ \\
 & {\scriptsize 67.4--68.2} & {\scriptsize 341.0--342.8} & {\scriptsize 1{,}003.8--1{,}009.2} & \\
\bottomrule
\end{tabular}
\end{table}

\subsection{The relay}

The relay was run end to end on a local Stellar network
(\file{stellar/quickstart}, core and RPC v29, protocol 28) with the Dfns client
signing through the Dfns mock. It confirmed an allowed transfer and an allowed
call whose tree nests an allowed transfer; it never had the custody key sign a
transfer over its cap or a call whose nested transfer the grant denies; it
refused a tampered entry at simulation without using a sequence number; and
it served two requests at once on two channels\claim{relay-local-network}.
Nido's unchanged TypeScript client relayed to confirmation through
it\claim{relay-frontend}. These are functional tests; the relay has no cost,
throughput or attack-suite measurement.

\subsection{Custody signing over channel accounts}

On Stellar testnet (protocol 29), three fresh accounts were given the custody
key as their only weighted signer (master weight 0, thresholds 1). For both
Dfns payload kinds that sign a Stellar transaction, \code{Hash} and
\code{Transaction}, a classic payment, a Soroban \code{InvokeHostFunction}
and a fee bump of a Soroban call, each sourced from a channel and signed only
by the custody key, were accepted; Horizon's records show every signature
verifying under the custody key. The same payment signed by the channel's
zeroed master key or by an unrelated key, a Soroban transfer signed by an
unrelated key, and a fee bump whose outer signature was the fee source's
zeroed master key were all refused with \code{txBAD\_AUTH} and used no
sequence number\claim{testnet-channel}. The custody key was held by the Dfns
mock, so this shows what the network accepts, not what Dfns will
sign (\S\ref{sec:relay}).

\paragraph{Not measured.} The local Stellar network for the grant signer and
the attack suite; the relay's cost, throughput and attack suite; throughput
against the size of the channel-account swarm; sibling spends submitted at
once through the swarm; and the host's signature authorizing a call through
a grant signer on chain. For concurrent sibling spends the protocol gives an
argument, not a measurement: two transactions whose footprints write the same
ledger entry conflict, and the result must be ``as if'' they were applied in
sequence~\cite{cap63}, so the later one in the protocol's order reads the
earlier one's charge to their common ancestors.

\section{Limitations and Open Questions}\label{sec:limits}

\paragraph{Access, not information flow.} A grant decides which actions an
agent may take. It does not track what the agent learns from an allowed action
or where that data goes next through other allowed actions: an agent allowed
to read a file and to reach a host can send the file's contents to that host.
Access control does not by itself enforce information-flow
policies~\cite{sabelfeld}; systems that track data flow through an agent's
tool calls~\cite{camel,fides} address a different question.

\paragraph{What the proofs cover.} The proofs are \translationtrust. The step
bound is proven for a counted model and is loose (\S\ref{sec:proofs}); the
measured cost follows the model's step count on the measured points, not
provably on every input\claim{deployed-cost-general}. The run-level budget
bound is proven for the grant signer's calls as \file{Signer.lean} models
them; that the contract's code is that model is
trusted\claim{wiring-budget-calls}.

\paragraph{The wiring is trusted, and has known gaps.} That the argument
mapping keeps distinct arguments distinct\claim{xdr-injective} and that the
relay's supplied and recorded trees cover every authorized
context\claim{trees-cover} are tested or argued, not proven. The host's path
resolution, URL parsing, redirect handling and payload decoding, and the
operating system calls beneath them, are tested only. Table~\ref{tab:gaps}
lists the gaps the paper states: under an empty grant the grant signer's own
address can authorize any call on another contract; a holder's signature
does not cover the grant id, so a key should hold one grant; no caveat binds
a signed transaction's fee, sequence number, time bounds or memo; network
caveats bind names, not the addresses they resolve to, and not the HTTP
method or path; and a file renamed after its descriptor was checked, or
hard-linked into the grant, is reached. The account's document must meet two
requirements nothing on chain checks\claim{document-requirements}; the grant
signer cannot check that the document names it, since it cannot see which
account delegated to it.

\paragraph{Scope of the hidden-call guarantee.} The grant signer sees the
contexts that need the account's authorization. Effects that need none, such
as a \code{transfer\_from} against a standing allowance or funds a contract
already holds for the account, are invisible to every caveat and budget, and
the paper assumes no allowed contract has such a
position\claim{no-standing-allowance}. The relay likewise cannot see calls
that ask for no authorization or calls a contract makes inside its own
\code{\_\_check\_auth}, and bounds the fee by a cap rather than a caveat.

\paragraph{Expressiveness.} Caveats are a closed set: membership of a field
in a list, an upper bound on one \code{i128} argument, and equality of one
argument. Each applies to every context of a tree alike, so a grant over two
contracts and two functions allows their cross product, and an argument
caveat reads the same position of every context (\S\ref{sec:grants}). There
is no caveat scoped to one contract and function, no lower bound, no expiry
and no caveat across two arguments. Allowances are cumulative and never reset:
the periodic and refillable allowances of session-key systems and Zodiac
Roles~\cite{zodiacRoles,smartSessions} have no counterpart, and would need a
time model in the budget code and its proofs. Budgets are enforced only at the
Soroban boundary, and each root grant meters one asset; transfers of other
assets are bounded only by caveats.

\paragraph{Keys and storage.} Holder keys at the Soroban boundary are ed25519
only. Grants live in persistent contract storage, extended to the network
maximum on every write; an archived grant fails authorization until it is
restored, which fails closed.

\paragraph{Evaluation.} The attack suite ran in Soroban simulation and at the
host, not on a local network and not through the relay; the relay has
functional tests only, with no cost, throughput or attack-suite measurement.
Host timings are the medians of one run on one machine.

\paragraph{Dfns.} Whether real Dfns's \code{Transaction} kind signs a
transaction whose source is a channel account\claim{dfns-transaction-source},
and whether its \code{Hash} kind signs 32 bytes directly as an ed25519
message\claim{dfns-hash}, are not stated by its reference and were not tried
against the service. The relay signs through \code{Hash}, which the testnet
run shows the network accepts when the signature is right.

\section{Related Work}\label{sec:related}

\paragraph{Capabilities and attenuation.} Dennis and Van Horn describe a
computation's authority as a list of capabilities, where a process ``may grant
some of its capabilities to the inferior sphere'' but cannot pass one ``not
implied by a capability present in the higher sphere''~\cite{dennis}. Miller's
principle of least authority and attenuation by interposition~\cite{miller},
SPKI's intersection of authorization tags along a certificate
chain~\cite{spki} and Capsicum's rule that a derived capability's rights
``must be a subset of the original rights''~\cite{capsicum} are the same
monotonicity that Theorems~\ref{thm:atten} and~\ref{thm:chain} state for
grants. Machine-checked versions of it are older and deeper than ours: seL4's
integrity proof covers its capability-based access control and holds ``via
refinement for the C implementation of the kernel''~\cite{sel4Integrity}, and
CHERI's proofs establish ``a reachable-capability monotonicity property'' of
the instruction set~\cite{cheriProofs}. Our monotonicity theorems are
immediate in the model; what is new is enforcing one verified evaluator's
verdicts at an agent host and inside a smart account's authorization. The
relay is a deputy holding the custody key for many agents; it avoids the
confused-deputy problem~\cite{confusedDeputy} by checking the requesting
agent's grant, not its own authority.

\paragraph{Caveat tokens.} Macaroons embed caveats, ``predicates that restrict
the macaroon's authority'', and let any holder derive a narrower credential by
adding one, with chained HMACs giving ``the inability for adversaries to remove
caveats''~\cite{macaroons}. Biscuit expresses checks in Datalog and uses
public-key signatures, so a holder ``cannot remove existing blocks without
invalidating the signature''~\cite{biscuit}. UCAN requires that ``each direct
delegation MUST either directly restate or attenuate (diminish) its
capabilities''~\cite{ucan}, and its delegation policies constrain an
invocation's arguments while leaving stateful conditions ``out of scope of the
UCAN Policy language''~\cite{ucanDelegation}. ZCAP capabilities ``inherit the
restrictions from all caveat properties of their parents''~\cite{zcapld}.
These are bearer credentials whose attenuation is enforced cryptographically;
what a caveat means is left to the verifier's code. Grants are held by the
enforcer and used by a holder key, the evaluator that gives caveats their
meaning is proven to compute the model, and budgets add the stateful
condition these token languages leave out: spending charged up the delegation
tree.

\paragraph{Smart-account delegation and session keys.} Account
abstraction lets ``each account to supply its own validation
logic''~\cite{erc4337}, and modular accounts add validator and hook
modules~\cite{erc7579}. ERC-7710 standardizes redeeming a delegation through
a delegation manager, leaving ``the structure of the
\code{\_permissionContexts}'' to each manager and recommending audits; it
names granting ``bounded permissions to AI agents'' as a
motivation~\cite{erc7710}. ERC-7715 lets an application request permissions
from a wallet, which ``MAY attenuate(reduce or increase)'' them~\cite{erc7715}.
MetaMask's delegation framework enforces caveats with enforcer contracts whose
hooks run ``from leaf to root delegation'' on redemption; a re-delegation can
add caveats, and ``to understand what the delegation can do, it is necessary to
analyze all the enforcers being used in the entire delegation
chain''~\cite{metamaskDelegation}. Its enforcers include per-delegation and
periodic transfer limits, and it is assured by audits; we found no formal
verification. Smart Sessions attaches policies with argument conditions and
ERC-20 spending limits to session keys~\cite{smartSessions}, and Zodiac Roles
evaluates condition trees over calldata with refillable
allowances~\cite{zodiacRoles}. On Stellar, OpenZeppelin's smart accounts check
context rules with policy contracts whose \code{enforce} ``both validates
conditions and applies state changes''~\cite{ozStellarAccounts}, and an agent
wallet example caps an agent's spending per period on one
token~\cite{ozAgentWallet}. Google's Agent Payments Protocol binds an agent's
payment to the user's intent with signed mandates, whose constraints ``allow
the verifier to verify that the Checkout and Payment match the User's
intent''~\cite{ap2}. These systems express richer conditions than our closed
caveat set, and several have the periodic or refillable budgets we lack. In
each, a caveat's meaning is the code of its enforcer or policy contract,
assured by audit and analysed chain by chain; here one evaluator gives every
caveat its meaning at every boundary, and attenuation, chaining and budget
safety are theorems about that evaluator's translation. The grant signer also
relies on Soroban handing a delegated signer every context of the account's
authorization~\cite{cap7101}, which a policy attached to one rule does not
see.

\paragraph{Custody policy engines.} Dfns evaluates amount and recipient rules
on plain value transfers only~\cite{dfnsPolicies}. Other engines decode
contract calls. Fordefi's policy conditions can specify ``the method's
parameters and their values''~\cite{fordefiPolicies}, with byte-string, array
and tuple parameters since May 2026~\cite{fordefiChangelog}; Turnkey's policy
language reads ``the parsed Ethereum transaction payload'' and, once a
contract's interface is uploaded, its function calls and
arguments~\cite{turnkeyPolicies,turnkeyContracts}; and Fireblocks' policy rules
match contract methods~\cite{fireblocksPolicyApi}, and its own comparison
claims ``method-level allow and deny by ABI function'' and transaction
simulation~\cite{fireblocksPolicy}. The engines whose documentation we read
describe rules over EVM-style calls; we found none that evaluates a Soroban
authorization tree. Grants differ from all four in three ways: holders
delegate by attenuation, the evaluator is verified, and the same policy is
enforced again on chain by the account itself, so the custody key is not the
only check.

\paragraph{Verified policy evaluation and Rust verification.} Cedar's authors
proved properties of a Lean model of the language and ``used extensive
differential random testing to ensure that the implementation matches the
model''~\cite{cedar}; they describe formally verifying ``the equivalence of the
deployed code and model'' as the alternative they did not take~\cite{cedarHow}.
Zelkova encodes AWS access policies into SMT to analyse them~\cite{zelkova}.
Our policy language is far smaller than Cedar's, and in exchange the deployed
Rust's translation is proven equal to the model through Charon and
Aeneas~\cite{charon,aeneas}. Rust can also be verified in place: Verus with
SMT and specifications written in Rust~\cite{verus}, Creusot through
Why3~\cite{creusot}, Prusti through Viper~\cite{prusti}, and hax by
translating to several provers' languages~\cite{hax}. We chose a functional
translation to Lean for two reasons: the readable model and the translation
live in one language, so the equivalence is an ordinary Lean theorem, and the
proofs use Lean's library and kernel rather than an SMT solver's.

\paragraph{WebAssembly sandboxing.} WebAssembly was designed with a formal
semantics~\cite{wasm}, mechanized in Isabelle~\cite{wasmcert}. vWasm and rWasm
compile Wasm into provably sandboxed code~\cite{bosamiya}, RLBox sandboxes
libraries inside Firefox with Wasm~\cite{rlbox}, and WaVe is a verified WASI
runtime ensuring that each file a sandbox accesses lies within its root
directory and that sockets reach only allow-listed addresses~\cite{wave}.
These isolate a module from resources it was not given. The agent host instead
gives an agent's tool calls a per-call check against a delegable grant, and
runs the checking evaluator itself as a component that imports
nothing~\cite{componentModelDocs}; it trusts Wasmtime~\cite{wasmtime} rather
than verifying the runtime, and confines no agent.

\paragraph{Agent security.} CaMeL tracks data flow through an agent's program
and enforces capability policies, written as Python functions, when tools are
called~\cite{camel}. FIDES ``propagates labels in messages, actions, tool calls
and results'' and executes consequential actions only when a policy over those
labels allows~\cite{fides}; both address the information flow that grants
leave out. Progent expresses privilege as ``symbolic rules over tool names and
arguments'' in JSON Schema and uses an SMT solver to decide whether a policy
update narrows or expands the agent's actions~\cite{progent}, and AgentSpec
lets users write rules with ``triggers, predicates, and enforcement
mechanisms'' evaluated at runtime~\cite{agentspec}; CaMeL and Progent evaluate
on AgentDojo~\cite{agentdojo}. A theory of what such gates can enforce shows
that ``deterministic pre-execution gates enforce exactly the nonempty
register-representable safety policies''~\cite{whatCanBeEnforced}; a grant
check is such a gate. nono
confines agent processes at the kernel with path, process and network rules
down to the HTTP method and path~\cite{nono}. Grants share Progent's monotonic
narrowing and nono's placement in front of the agent's resources, but take
their checks down to argument values, carry them past the agent's process into
the smart account that executes the transaction, and prove the evaluator that
decides them.

\section{Conclusion}\label{sec:conclusion}

An agent's authority can be one policy enforced at two boundaries. Written as
a grant, it narrows with every delegation and carries a budget charged up the
delegation tree, and the evaluator that decides it computes a readable model,
\translationtrust. In the Soroban smart account the evaluator sees every
context of the account's authorization through CAP-0071 delegated signing, at
2.2M to 5.0M of a transaction's 400M instructions, and refuses every one of
the attack suite's 45 attempts in simulation without changing state. On a
local Stellar network, with every call sent through the relay, the chain
refuses the 44 that run there whether or not the relay's own grant stops them
first, and sub-agents racing to overdraw their root are charged within its
allowance. In the agent host the evaluator decides file, network and signing
requests in about a microsecond per check. Its metered work follows the
model's step count on 27 of the 28 measured pairs; on the 28th, where both
checks scan a whole list, the two disagree in sign by a few steps. What remains trusted is the translation, the compilers and the wiring that
turns calls into actions, including the grant signer's storage taken as the
Lean model of its budget path; the paper states the wiring's remaining gaps,
among them that a holder's signature does not cover the grant id and that no
caveat binds a signed transaction's fee.


\bibliographystyle{ACM-Reference-Format}
\bibliography{references}

\appendix
\onecolumn
\section{Claim Ledger}\label{app:claims}

Every claim in this paper and its slides, generated from
\file{paper/claims.yaml}. A claim is \claimtag{proven} by a Lean theorem,
\claimtag{tested} by a named test, \claimtag{measured} by a recorded
measurement, \claimtag{assumed} as a stated assumption, or
\claimtag{unproven}. The build checks that every theorem is declared in its
file, every test name appears in its file, every measurement file exists, and
every claim is used by the paper or the slides.

{\footnotesize
\begin{xltabular}{\linewidth}{@{}>{\raggedright\arraybackslash}p{0.15\linewidth}
    >{\raggedright\arraybackslash}X l
    >{\raggedright\arraybackslash}p{0.30\linewidth}@{}}
  \caption{All claims and their status.}\label{tab:claims}\\
  \toprule
  \textbf{Claim} & \textbf{Statement} & \textbf{Status} & \textbf{Evidence}\tabularnewline
  \midrule
  \endfirsthead
  \toprule
  \textbf{Claim} & \textbf{Statement} & \textbf{Status} & \textbf{Evidence}\tabularnewline
  \midrule
  \endhead
  \bottomrule
  \endlastfoot
  \ClaimTableRows
\end{xltabular}}


\end{document}
